\chapter{Reproducibilidad y archivos técnicos}
\section{Manifiesto mínimo de un ensayo}
\begin{enumerate}
\item Fecha, identificador, caso, variante y propósito del ensayo.
\item Origen y huella de datos, modelo, pesos y código fuente.
\item Versión del intérprete, bibliotecas, convertidor, compilador y entorno de ejecución.
\item Opciones de generación, optimización, frecuencias y asignación de memorias.
\item Protocolo de medición, calentamiento, repeticiones y frontera energética.
\item Registros sin procesar y procedimiento que genera tablas y figuras.
\end{enumerate}

\section{Verificación del binario}
Para conservar el desensamblado se puede utilizar la herramienta de la cadena Arm:
\begin{quote}\ttfamily\small
arm-none-eabi-objdump -d -S aplicacion.elf > disassembly.txt
\end{quote}
La búsqueda de VMLADAV(A), instrucciones de carga vectorial y predicación permite localizar candidatos MVE. El nombre completo de función, su llamada y el perfil asociado deben quedar registrados. La búsqueda textual por sí sola no prueba ejecución de una instrucción ni mide su contribución a la latencia.

\section{Organización sugerida de evidencias}
\begin{table}[H]\centering\small
\caption{Archivos que acompañan las mediciones.}
\begin{tabularx}{\linewidth}{p{4cm}Y}\toprule
Directorio & Contenido\\\midrule
\texttt{modelos/} & Modelos por formato y SHA-256\\
\texttt{datos/} & Particiones y manifiestos, según licencia\\
\texttt{reportes\_st/} & Análisis, perfil y asignación CPU/NPU\\
\texttt{firmware/} & Fuentes, ELF, mapa y binarios\\
\texttt{registros/} & Salidas UART y mediciones sin procesar\\
\texttt{analisis/} & Scripts de métricas y figuras\\\bottomrule
\end{tabularx}\end{table}
\pendiente{Incorporar aquí los fragmentos esenciales del código validado. Los programas completos se entregan como material complementario y se referencian por versión.}
